\section{Priority Research Directions and Experimental Design}
\label{sec:future-directions}

The following schemes are proposed jointly from the existing evidence and the model candidates; all are experimental hypotheses rather than verified recipes for the target phase. We recommend proceeding in the order ``local phase diagram--site occupancy/oxygen stoichiometry--crystal growth--analogous substitution''; the phase-composition results of each stage determine whether the next stage is worth undertaking.

\subsection{Direction 1: A Local Phase Diagram near the Target Composition}

The first priority is to establish a fine-step composition grid around $\mathrm{Ba_5Y_{12}ZnSi_8O_{40}}$. The base precursors are the model's top-ranked and chemically conventional BaCO$_3$, Y$_2$O$_3$, ZnO, and SiO$_2$, with the target stoichiometry corresponding to a molar ratio of $5{:}6{:}1{:}8$. This direction is a model prediction awaiting experimental verification. We suggest first screening in stages over the temperature range of the Zn-free solid-state studies: an initial reaction at about 1000\,$^\circ$C for 12\,h, followed by regrinding and annealing at 1300\,$^\circ$C for 24\,h, with repetition decided on the basis of the PXRD results; short-duration points at 1450--1600\,$^\circ$C are added only if obvious unreacted material remains and the mass loss and crucible compatibility are controllable. The temperatures and times are taken from the screening ranges of the neighboring Ba--Y solid-state studies~\cite{yamane2024microstructure,gulay2024navigation}; they serve only to organize the experimental matrix and do not imply that the target phase necessarily forms under these conditions.

The composition axes should include at least: (i) small deviations of the Ba\slash Y ratio around the target point; (ii) variation of the Zn content; and (iii) the Zn-free reference phase related to the oxygen stoichiometry. For every sample, the pre- and post-firing masses, the thermal history, and the cooling mode are recorded, and the phase composition and measured element ratios are given by quantitative PXRD\slash Rietveld analysis, SEM--EDS\slash EPMA, and ICP analysis where necessary. If glass or a multiphase boundary appears, sampling should be densified in that region rather than retaining only the most ideal diffraction pattern.

\subsection{Direction 2: Zn--Y--Oxygen-Stoichiometry Coupling}

A testable charge-balanced series is
\[
\mathrm{Ba_5Y_{13-\mathit{x}}Zn_{\mathit{x}}[SiO_4]_8O_{8.5-\mathit{x}/2}},
\qquad 0\leq x\leq1.25,
\]
where $x=0$ corresponds to the reported Zn-free tetragonal reference phase and $x=1$ corresponds, in elemental stoichiometry, to the target composition. We suggest $x=0,0.25,0.50,0.75,1.00,1.25$ as the discrete points of the first round. This direction is a model prediction awaiting experimental verification. The precursors remain BaCO$_3$, Y$_2$O$_3$, and SiO$_2$, with the Zn content adjusted by ZnO; the expression is only a charge-balance hypothesis and does not presuppose the true site occupancy of Zn or the distribution of oxygen vacancies. Assessment should examine simultaneously the continuity of the lattice parameters, the phase fractions, the measured Zn\slash Y ratio, and the site-occupancy refinement; where feasible, neutron diffraction or independent measurements of oxygen content and valence state can further distinguish Zn substitution, oxygen defects, and a two-phase mixture.

The scientific value of this series exceeds that of a single-point reproduction: a continuous change in the structural parameters supports limited solid solution; essentially unchanged lattice parameters with a changing two-phase ratio is closer to a two-phase region; and new diffraction peaks appearing only at $x\approx1$ suggest a line compound. All three outcomes constrain the formation mechanism of the target phase.

\subsection{Direction 3: Solid-State Reaction in Parallel with High-Temperature Solution Crystal Growth}

The solid-state reaction route serves to locate the bulk phase region, and the high-temperature solution route serves to obtain crystals suitable for SCXRD analysis. This direction is a model prediction awaiting experimental verification. The high-temperature solution route can take the K$_2$CO$_3$\slash MoO$_3$ flux system already validated in the Ba--Y--Si--O system as its reference starting point, with the BaCO$_3$, Y$_2$O$_3$, ZnO, and SiO$_2$ precursors of Direction 1; the maximum temperature is set within 1100--1200\,$^\circ$C, held for about 3\,h, followed by slow cooling at 1--2\,K\,h$^{-1}$ to about 900\,$^\circ$C, while the nutrient-to-flux ratio and the Y$_2$O$_3$\slash SiO$_2$ ratio are varied systematically~\cite{kolitsch2009crystal,wierzbickawieczorek2017high}. These values come from crystal growth in the Zn-free homologous series; they must be compared in parallel with the solid-state reaction route using the target starting composition and cannot be called a recipe for the target phase.

Each crystal-growth batch should report at least the liquid-phase composition, the crucible filling ratio, the maximum temperature, the dwell time, the cooling rate, the separation method, the crystal size, and the yield; the crystals are verified by SCXRD and EDS, and the residual mother liquor is checked by PXRD and compositional analysis. If only a few target single crystals are obtained while the composition of the mother liquor is far from the target composition, this should be interpreted as selective growth from the liquid phase rather than as a bulk phase region.

\subsection{Direction 4: Mg\slash Co Analogs and Model-Proposed Precursors}

This review invoked a local two-step inorganic retrosynthesis model to predict Top-5 precursor combinations for three elemental formulae. The model's top-ranked combinations for the Zn, Mg, and Co targets all retain BaCO$_3$, Y$_2$O$_3$, and SiO$_2$, with ZnO, MgO, and Co$_3$O$_4$ as the fourth starting material, respectively; the probabilities of the three top-ranked combinations within their respective candidate pools are 0.594, 0.522, and 0.298. This direction is a model prediction awaiting experimental verification. Table~\ref{tab:mcp-experiments} converts the model output into minimal testable experiments. The model proposes precursor combinations solely on the basis of elemental stoichiometry and is not used to judge the silicate structural units; the text throughout retains the structural formula used in the target publication.

\begin{table}[bp]
  \centering
  \caption{\textbf{Top-ranked precursor combinations from the two-step inorganic retrosynthesis model and their experimental uses.} All combinations are model-ranked results (not validated by chemical routes); the process conditions are taken from the literature on neighboring systems and require re-checking by phase-diagram experiments.}
  \label{tab:mcp-experiments}
  \footnotesize
  \begin{tabular}{P{\dimexpr 0.1188\textwidth-2\tabcolsep\relax}P{\dimexpr 0.3100\textwidth-2\tabcolsep\relax}P{\dimexpr 0.1300\textwidth-2\tabcolsep\relax}P{\dimexpr 0.4412\textwidth-2\tabcolsep\relax}}
    \toprule
    \thead{Hypothe\-tical target} & \thead{Model top-ranked precursors} & \thead{Candidate-pool probability} & \thead{Suggested experiment and criteria}\\
    \midrule
    Zn target & BaCO$_3$ + Y$_2$O$_3$ + SiO$_2$ + ZnO & 0.594 & Base starting composition for the local phase diagram; solid-state screening first, then transfer of near-single-phase compositions to flux growth; bulk phase region and structural assignment determined by PXRD\slash Rietveld analysis and SCXRD, respectively\\
    \addlinespace[3pt]
    Mg analog & BaCO$_3$ + Y$_2$O$_3$ + SiO$_2$ + MgO & 0.522 & Mg$^{2+}$ as an isovalent size control; screening along the same composition/thermal-history matrix as for Zn; substitution on the same site discussed only if a new phase appears and the measured composition varies continuously\\
    \addlinespace[3pt]
    Co analog & BaCO$_3$ + Y$_2$O$_3$ + SiO$_2$ + Co$_3$O$_4$ & 0.298 & Air atmosphere as the first-round condition; measured composition and phase fractions determined; secondary phases not to be mistaken for structural substitution\\
    \bottomrule
  \end{tabular}
\end{table}

The remaining Zn routes given by the model include additional BaO, BaF$_2$, or Ba(NO$_3$)$_2$, or non-standard database labels; their probabilities are all clearly lower than that of the top-ranked combination, and they should not be recommended alongside it. An additional Ba source would alter the local phase diagram; BaF$_2$ increases the risk of fluorine-containing volatiles and waste streams; and non-standard labels may further reflect normalization problems in the training data. MgCO$_3$ in the Mg routes can serve as a separate control for precursor form but should not be added together with MgO; the several Co oxide precursors should not be regarded as three equivalent schemes.

\subsection{Organizing Experiments for Scientific Discovery}

Each round of experiments should be judged by the criterion of ``maximizing distinguishable conclusions'' rather than aiming merely at obtaining crystals. The most informative combination is: (i) a fine-step composition grid coupled with quantitative phase analysis to establish the local phase diagram; (ii) parallel solid-state reaction and flux-growth experiments at the same composition to distinguish the equilibrium phase region from selective nucleation in the liquid phase; and (iii) a coupled Zn-content and oxygen-stoichiometry series to test the site and defect hypotheses. All negative results are likewise preserved, because they determine the composition choices of the next round and constrain the boundary of ``synthesizability''.
